% Generated by uscribe -- do not edit by hand
\documentclass[letterpaper,12pt]{article}
\usepackage[T1]{fontenc}
\usepackage[utf8]{inputenc}
\usepackage{graphicx}
\usepackage{hyperref}
\usepackage{xcolor}
\usepackage{fancyvrb}
\usepackage{booktabs}
\usepackage{longtable}
\IfFileExists{multirow.sty}{\usepackage{multirow}}{}
\hypersetup{colorlinks=true,linkcolor=blue,urlcolor=blue}
\begin{document}
\begin{titlepage}
\null\vfil\vskip 40pt
\begin{center}
{\LARGE\bfseries Multicast and Socket Attributes in Unicon \par}
\vskip 3em
{\large\bfseries Jafar Al-Gharaibeh \par}
\vskip 1.0em
{\bfseries Unicon Technical Report: 28\par}
\vskip 1.0em
{\large\bfseries August 2026\par}
\vskip 5.5em
{\bfseries Abstract}\par
\begin{quote}
Unicon could open TCP and UDP sockets, but applications had no way to join a multicast group, pick an interface, or set TTL from the language. This report describes trailing name=value attributes on open(), the same style as SSL and graphics windows, plus Attrib() after open() for mutation and f["name"] / key(f) for status. Binding a UDP socket to a group address joins that group; source@group or source= selects source-specific multicast. Defaults cover reuse, all-interface IPv4 joins, and outbound interface selection so the common cases need no attributes at all.
\end{quote}
\begin{quote}
{\bfseries Keywords:} Unicon, multicast, SSM, UDP, socket attributes, Attrib, runtime, technical report.
\end{quote}
\vskip 0.8in
{\large\bfseries
Unicon Project\\
\url{http://unicon.org}\\
\ \
University of Idaho\\
Department of Computer Science\\
Moscow, ID, 83844, USA
}
\end{center}
\vfil\null
\end{titlepage}
\setcounter{page}{1}
\section{1. Introduction}
\label{utr28-sec-intro}

This report describes IP multicast [7] and the socket-attribute facility it sits on. Trailing arguments to \texttt{open()} on network modes are \texttt{name=value} strings applied with \texttt{setsockopt(2)}. After \texttt{open()}, \texttt{Attrib()} uses the same names. The language-facing reference lives in \emph{Programming with Unicon} [8] (Chapter 5, "Multicast and Socket Attributes") and the language reference (\texttt{open} trailing attributes, \texttt{Attrib}, subscript peek, \texttt{key()}). Automated tests live in \texttt{tests/posix/mcast.icn} and \texttt{tests/posix/sockpeek.icn}.


The work is part of POSIX networking, not an optional library. A later report covers raw sockets [9], which reuse these attributes (\texttt{ttl}, \texttt{iface}, \texttt{join}) on mode \texttt{nr}.

\section{2. Motivation}
\label{utr28-sec-motivation}

Unicon's \texttt{open()} already constructed TCP and UDP handles. Multicast -- any-source (ASM) and source-specific (SSM) [10] -- is ordinary UDP plus membership and hop-limit options. Without those options, a Unicon program could not join \texttt{239.1.1.1}, restrict a join to one NIC, or send with a limited TTL. Those operations were not supported.


The rest of the I/O API already has the pattern: \texttt{open()} constructs, trailing strings configure (SSL \texttt{key=}, window \texttt{WAttrib}), \texttt{Attrib()} mutates later, and \texttt{f["name"]} peeks status. Socket options follow that pattern instead of adding \texttt{setsockopt()} as a new global.

\section{3. Design principles}
\label{utr28-sec-principles}

\textbf{Attributes are \texttt{name=value}, like SSL.} Unknown names fail \texttt{open()} with error 1310 (\texttt{bad~socket~attribute}). Booleans are exactly \texttt{yes} or \texttt{no}, matching \texttt{verifyPeer}.

\begin{Verbatim}[fontsize=\small,frame=single,commandchars=\\\{\}]
procedure main()
   if open(":5110", "nua4", "bogus=yes") then
      write("oops")
   else
      write("bogus: ", &errortext)
   if open(":5110", "nua4", "reuseaddr=1") then
      write("oops")
   else
      write("numeric boolean: ", &errortext)
end
\end{Verbatim}


Output:

\begin{Verbatim}[fontsize=\small,frame=single,commandchars=\\\{\}]
bogus: bad socket attribute
numeric boolean: bad socket attribute
\end{Verbatim}


\textbf{DWIM for the common cases.} Binding a UDP socket to a multicast group joins that group. Listeners default to \texttt{reuseaddr=yes}. No \texttt{iface=} means join every UP IPv4 interface and send out the first non-loopback IPv4 address (else loopback). Sending to \texttt{255.255.255.255} enables \texttt{SO\_BROADCAST}. Explicit attributes always override.


\textbf{\texttt{source@group}, not \texttt{group@source}.} The address form matches VLC, ffmpeg, and RFC 4607 \texttt{(S,G)} order. Empty either side is invalid.


\textbf{\texttt{ttl} and \texttt{iface}, not multicast-only names.} Hop limit and interface are used by unicast, multicast, and later raw sockets [9]. \texttt{ttl=} sets both \texttt{IP\_TTL} and \texttt{IP\_MULTICAST\_TTL} (or the IPv6 hop options). \texttt{iface=} accepts an IPv4 address, a numeric index, or a name (\texttt{en0}, \texttt{lo0}; Windows also maps \texttt{lo} / \texttt{lo0} to \texttt{loopback\_0}).


\textbf{Order matters.} \texttt{iface=} before \texttt{join=} selects the interface used by that join. \texttt{Attrib()} applies a batch of assignments together for the same reason.


\textbf{Cached listeners are aliases; attributes get a new socket.} A second \texttt{open()} of the same listener address with no socket attributes returns another file for the cached fd. An open that supplies \texttt{join=} / \texttt{iface=} / and so on creates an independent socket so it cannot retune earlier aliases. Use \texttt{Attrib()} to change options on a live shared handle.

\section{4. Opening a multicast socket}
\label{utr28-sec-open}

IP multicast uses UDP (modes \texttt{nu} / \texttt{nua}, optionally with \texttt{4} or \texttt{6} for family). A minimal receiver and sender:

\begin{Verbatim}[fontsize=\small,frame=single,commandchars=\\\{\}]
procedure main()
   local f, r
   f := open("239.1.1.1:5000", "nua") | stop(&errortext)
   while r := receive(f) do
      write(r.addr, "  ", r.msg)
end
\end{Verbatim}

\begin{Verbatim}[fontsize=\small,frame=single,commandchars=\\\{\}]
procedure main()
   local f, i
   f := open("239.1.1.1:5000", "nu") | stop(&errortext)
   every i := 1 to 5 do
      writes(f, "hello-" || i) & delay(500)
end
\end{Verbatim}


The receiver binds the group address, so the runtime joins that group (ASM) and enables port reuse. The sender's destination is the group; without \texttt{iface=}, outbound multicast uses the first non-loopback IPv4 address. Kernel multicast loopback stays on, so a local receiver still sees the packets.


The same exchange in one process over loopback (this is the shape of \texttt{tests/posix/mcast.icn}). \texttt{iface=127.0.0.1} keeps the datagram on loopback even when the host has no multicast route:

\begin{Verbatim}[fontsize=\small,frame=single,commandchars=\\\{\}]
procedure main()
   local f, s, r
   f := open("239.42.42.99:5199", "nua", "iface=127.0.0.1") |
        stop(&errortext)
   s := open("239.42.42.99:5199", "nu", "iface=127.0.0.1") |
        stop(&errortext)
   writes(s, "hello-multicast")
   if *select(f, 2000) > 0 then \{
      r := receive(f)
      write(r.msg)
      \}
   else
      write("timeout")
   close(s)
   close(f)
end
\end{Verbatim}


Output:

\begin{Verbatim}[fontsize=\small,frame=single,commandchars=\\\{\}]
hello-multicast
\end{Verbatim}


\texttt{r.addr} is the sender (host plus an ephemeral port). Alpine often does not deliver loopback multicast even when the join succeeds; the test suite treats a missing datagram as success there so CI stays portable.


A wildcard bind plus explicit joins is the other shape:

\begin{Verbatim}[fontsize=\small,frame=single,commandchars=\\\{\}]
f := open(":5000", "nua",
          "iface=192.168.1.10",
          "join=239.1.1.1",
          "join=232.1.1.1,10.0.0.5")
\end{Verbatim}


\texttt{join=group} is ASM; \texttt{join=group,source} is SSM. Attributes may be repeated. Mixing ASM and SSM in one call is allowed. If \texttt{open()} has no \texttt{4} / \texttt{6} flag, a \texttt{join=} group address picks the family so an IPv4 join is not attempted on an IPv6 wildcard socket.

\section{5. Source-specific multicast}
\label{utr28-sec-ssm}

Name the source before the group, or pass \texttt{source=} on a socket already bound to the group. Both IPv4 and IPv6 are supported. For IPv6 SSM, pass an explicit \texttt{iface=} (name or index) so the join uses a single interface:

\begin{Verbatim}[fontsize=\small,frame=single,commandchars=\\\{\}]
f := open("192.168.1.1@239.1.1.1:5000", "nua")
f := open("239.1.1.1:5000", "nua", "source=192.168.1.1")
f := open("::1@ff3e::1:5000", "nua6", "iface=lo0")
f := open("ff3e::1:5000", "nua6", "iface=lo0", "source=::1")
\end{Verbatim}


On send or connect, \texttt{source@group:port} keeps only the group so the destination is the group address. The source is a receive-side filter, not a send address.


IPv6 SSM is a single \texttt{MCAST\_JOIN\_SOURCE\_GROUP} / \texttt{MCAST\_LEAVE\_SOURCE\_GROUP} with no multi-interface walk. Those walks have triggered macOS kernel panics in IPv6 source filter teardown. Already-a-member and already-gone are treated as success.

\section{6. Attributes}
\label{utr28-sec-attrs}
\begin{longtable}{ll}
\hline
Attribute & Role \\
\hline
\endfirsthead
\hline
Attribute & Role \\
\hline
\endhead
\hline
\endfoot
\hline
\endlastfoot
\texttt{join=} & Add membership: \texttt{group} (ASM) or \texttt{group,source} \\
\texttt{leave=} & Drop membership; same forms as \texttt{join}. Pass \\
\texttt{source=} & SSM source for a group bind (repeatable). Set-only. \\
\texttt{iface=} & IPv4 address, index, or name. Restricts join and \\
\texttt{ttl=} & Unicast and multicast hop limits. \texttt{f["ttl"]} peeks \\
\texttt{mcastloop=yes|no} & Whether the host receives its own multicasts. Kernel \\
\texttt{reuseaddr=} / \texttt{reuseport=} & Bind reuse. Listeners default to \texttt{reuseaddr=yes} \\
\texttt{broadcast=} & \texttt{SO\_BROADCAST}. Opening \texttt{255.255.255.255} sets it. \\
\texttt{rcvbuf=} / \texttt{sndbuf=} & Buffer sizes in bytes. \\
\end{longtable}


\texttt{proto} and \texttt{hdrincl} are listed in the same table in the runtime; they belong to raw sockets [9].

\section{7. \texttt{Attrib()} and status peek}
\label{utr28-sec-attrib}
\begin{Verbatim}[fontsize=\small,frame=single,commandchars=\\\{\}]
procedure main()
   local f, g
   f := open(":5110", "nua4") | stop(&errortext)
   Attrib(f, "ttl=4")
   write("ttl=", f["ttl"])
   write("mcastloop=", f["mcastloop"])
   if g := f["groups"] then
      every write(!g)
   else
      write("groups: unpopulated")
   close(f)
end
\end{Verbatim}


Output:

\begin{Verbatim}[fontsize=\small,frame=single,commandchars=\\\{\}]
ttl=4
mcastloop=yes
groups: unpopulated
\end{Verbatim}


Assignments in one call are applied together, preserving \texttt{iface} then \texttt{join} order. Status is peeked with \texttt{f["name"]}; \texttt{Attrib()} only assigns. Unknown names raise

\begin{enumerate}
\item 
An unpopulated field fails. Boolean fields that answered
\end{enumerate}

succeed with \texttt{"yes"} or \texttt{\&null} (never \texttt{"no"}). \texttt{key(f)} generates every answerable field. \texttt{f["*"]} snapshots those fields under one lock. Peeking a closed handle is error 174. \texttt{key(f)} that has not yet produced a name also raises 174 if the handle is already closed. After the first suspend, \texttt{close()} makes the generator fail instead of raising, so a walk does not turn a mid-generation close into an error.


TCP, UDP, multicast, and raw sockets share one peek table (\texttt{sock\_peek}). \texttt{join}, \texttt{leave}, and \texttt{source} are verbs, not peek fields. \texttt{proto} is stored at \texttt{socket()} time on raw sockets [9]. \texttt{f["groups"]} is a list of joined groups (SSM as \texttt{source@group}); a single group is still a list of length 1.

\begin{longtable}{ll}
\hline
Name & Value \\
\hline
\endfirsthead
\hline
Name & Value \\
\hline
\endhead
\hline
\endfoot
\hline
\endlastfoot
\texttt{ttl} & Hop limit (integer), typically \texttt{1} .. \texttt{255} \\
\texttt{iface} & Dotted IPv4 address (\texttt{127.0.0.1}) or IPv6 ifindex (\texttt{1}) \\
\texttt{groups} & List of joined groups, e.g. \texttt{239.1.1.1}; SSM as \\
\texttt{mcastloop} & \texttt{"yes"} or \texttt{\&null} \\
\texttt{reuseaddr} / \texttt{reuseport} & \texttt{"yes"} or \texttt{\&null} (\texttt{reuseport} may be unpopulated) \\
\texttt{broadcast} & \texttt{"yes"} or \texttt{\&null} \\
\texttt{rcvbuf} / \texttt{sndbuf} & Buffer sizes in bytes (integers) \\
\texttt{proto} & IP protocol number: \texttt{1} (ICMP), \texttt{2} (IGMP), \texttt{89} \\
\texttt{hdrincl} & \texttt{"yes"} or \texttt{\&null}; typical on raw sockets \\
\end{longtable}

\begin{Verbatim}[fontsize=\small,frame=single,commandchars=\\\{\}]
f := open(":5110", "nua4") | stop(&errortext)
Attrib(f, "ttl=4")
Attrib(f, "iface=127.0.0.1", "join=239.1.1.1")
Attrib(f, "iface=127.0.0.1", "join=239.1.1.2")
write("ttl=", f["ttl"])
every g := !f["groups"] do
   write("joined ", g)
every k := key(f) do
   write("field ", k)
close(f)
\end{Verbatim}


Membership changes still use assignment form:

\begin{Verbatim}[fontsize=\small,frame=single,commandchars=\\\{\}]
Attrib(f, "join=239.1.1.2")
Attrib(f, "iface=127.0.0.1", "leave=239.1.1.1")
\end{Verbatim}


On some Linux/musl builds \texttt{getsockopt(IP\_MULTICAST\_IF)} does not recover the interface used at join time. Pass \texttt{iface=} on \texttt{leave=} so \texttt{IP\_DROP\_MEMBERSHIP} matches the join. Already-not-a-member is success; the runtime also retries \texttt{INADDR\_ANY} and each local IPv4 address when a single DROP fails.

\section{8. Listener cache}
\label{utr28-sec-cache}

\texttt{sock\_listen} still caches the first bind of an address so repeated \texttt{open()} of the same listener is cheap. That cache is a problem if a later \texttt{open(...,~"join=...")} reused the fd and changed memberships for every live alias.


Rules:

\begin{itemize}
\item 
Cache hit, no socket attributes: another \texttt{File} for the same fd (owner refcount). Closing one alias leaves the others until the last handle is closed.
\item 
Cache hit with attributes: release the pin and create an independent socket (\texttt{reuseaddr} makes the second bind succeed). \texttt{leave=} on that socket does not drop the cached alias's membership.
\item 
Changing options on a shared handle: \texttt{Attrib()} on one of the aliases.
\end{itemize}

A cache-hit open that is not a pure alias only allows additive \texttt{join=} / \texttt{source=}; \texttt{leave=} and option changes on that path would retune every live \texttt{File}.

\section{9. Runtime implementation}
\label{utr28-sec-impl}

All of this lives in \texttt{src/runtime/rposix.r}, applied from \texttt{open()} in \texttt{fsys.r} and from \texttt{Attrib()} in \texttt{fmisc.r}.


\textbf{Two passes.} \texttt{reuseaddr} / \texttt{reuseport} must be set between \texttt{socket()} and \texttt{bind()}. Everything else runs after the socket is bound. \texttt{apply\_sock\_attrs()} takes a \texttt{prebind} flag.


\textbf{Auto-join.} After the post-bind pass, a UDP socket bound to a multicast group is joined unless an explicit \texttt{join=} took over. \texttt{source@group} or \texttt{source=} makes that join SSM.


\textbf{IPv4 join without \texttt{iface=}.} Walk \texttt{getifaddrs()} and \texttt{IP\_ADD\_MEMBERSHIP} on every UP IPv4 address so same-host and multi-homed receives work without naming a NIC. An explicit \texttt{iface=} joins only that one. Outbound multicast without \texttt{iface=} picks the first non-loopback IPv4 address via \texttt{IP\_MULTICAST\_IF}.


\textbf{IPv6 ASM} uses \texttt{IPV6\_JOIN\_GROUP} / \texttt{IPV6\_LEAVE\_GROUP} with an ifindex (zero means default). There is no all-iface walk like IPv4.


\textbf{Errors.} Unknown attributes and bad booleans are 1310. \texttt{setsockopt} failures fail \texttt{open()} / \texttt{Attrib()} with the system \texttt{\&errortext}, except idempotent membership errors (already a member / already gone).


Mode \texttt{n} with \texttt{4} or \texttt{6} still selects the address family. \texttt{"nua4"} is useful when a join group must match an IPv4 wildcard.

\section{10. Status and remaining work}
\label{utr28-sec-status}

The language surface described here is implemented: trailing attributes, implicit ASM/SSM joins, \texttt{Attrib()} join/leave, listener-cache isolation, IPv4 all-iface join, IPv6 SSM with an explicit interface, and the broadcast shortcut.


\textbf{Membership query.} \texttt{f["groups"]} is a list of groups joined through \texttt{open()} or \texttt{Attrib()} \texttt{join=} / \texttt{source=}. The kernel still has no portable membership dump; this is runtime-tracked state.


\textbf{IPv6 all-interface join.} IPv4 walks every NIC; IPv6 does not. Callers who need every IPv6 interface must join per \texttt{iface=}.


\textbf{IPv6 SSM without \texttt{iface=}.} Not recommended. The implementation refuses to walk interfaces for IPv6 source filters.


\textbf{Loopback delivery.} Alpine (especially under QEMU) often does not deliver loopback multicast even when the join succeeds. \texttt{tests/posix/mcast.icn} still checks membership and treats a missing datagram as success on Alpine so CI stays portable.


\textbf{Windows IPv6 SSM.} Membership calls are attempted; \texttt{loopback\_0} and UNIX \texttt{lo} / \texttt{lo0} names are mapped. Delivery and option support vary by stack; the test skips cleanly on \texttt{ENOPROTOOPT}.


\textbf{Raw sockets} [9] reuse \texttt{ttl} / \texttt{iface} / \texttt{join} on \texttt{SOCK\_RAW}. That is a separate report.

\section{Appendix: test coverage}
\label{utr28-sec-coverage}

\texttt{tests/posix/mcast.icn} runs in one process over loopback, one UDP port per case so a late datagram cannot land on the next receiver. It checks:

\begin{itemize}
\item 
unknown attributes and numeric booleans fail with 1310
\item 
explicit \texttt{join=} plus \texttt{iface=127.0.0.1}
\item 
implicit join by binding the group address
\item 
SSM via \texttt{source@group} and via \texttt{source=}
\item 
broadcast open of \texttt{255.255.255.255} (never \texttt{EACCES})
\item 
\texttt{Attrib()} \texttt{ttl=}, \texttt{join=}, \texttt{leave=}
\item 
IPv6 SSM membership via \texttt{@}, \texttt{source=}, and \texttt{join=group,source} (delivery not required)
\item 
cache-hit alias vs independent attr open
\end{itemize}

\texttt{tests/posix/sockpeek.icn} checks \texttt{f["ttl"]}, \texttt{f["groups"]} as a list after two joins, and \texttt{key(f)}.


Expected output is \texttt{tests/posix/stand/mcast.std}:

\begin{Verbatim}[fontsize=\small,frame=single,commandchars=\\\{\}]
bogus attribute: bad socket attribute
numeric boolean: bad socket attribute
Received hello-multicast
Received hello-implicit
Received hello-ssm-at
Received hello-ssm-source
broadcast: ok
ttl=3
Received hello-attrib
leave: ok
ipv6-ssm-at: ok
ipv6-ssm-source: ok
ipv6-ssm-join: ok
cache-attr open: independent
Received hello-alias
\end{Verbatim}

\section{References}
\label{utr28-references}
\begin{thebibliography}{99}
\bibitem[1]{cite-jeffery-utr15}
Clinton Jeffery. "How to Write a Unicon Technical Report". doc/utr/utr15.tex. 2013.

\bibitem[2]{cite-libssh}
The libssh Project. "libssh -- The SSH Library". https://www.libssh.org/. 2026.

\bibitem[3]{cite-jeffery-utr13}
Clinton Jeffery. "The Unicon Messaging Facilities". doc/utr/utr13.odt.

\bibitem[4]{cite-openssl}
The OpenSSL Project. "OpenSSL -- Cryptography and SSL/TLS Toolkit". https://www.openssl.org/. 2026.

\bibitem[5]{cite-algharaibeh-utr26}
Jafar Al-Gharaibeh. "Native SSH and SFTP Support in Unicon". doc/utr/utr26.rst. 2026.

\bibitem[6]{cite-rfc2104}
H. Krawczyk and M. Bellare and R. Canetti. "HMAC: Keyed-Hashing for Message Authentication". RFC 2104. 1997.

\bibitem[7]{cite-rfc1112}
S. Deering. "Host Extensions for IP Multicasting". RFC 1112. 1989.

\bibitem[8]{cite-jeffery-pwu}
Clinton Jeffery and Shamim Mohamed and Jafar Al-Gharaibeh. "Programming with Unicon". doc/book/. 2026.

\bibitem[9]{cite-algharaibeh-utr29}
Jafar Al-Gharaibeh. "Raw Sockets and Packet Layouts in Unicon". doc/utr/utr29.rst. 2026.

\bibitem[10]{cite-rfc4607}
H. Holbrook and B. Cain. "Source-Specific Multicast for IP". RFC 4607. 2006.

\bibitem[11]{cite-algharaibeh-utr28}
Jafar Al-Gharaibeh. "Multicast and Socket Attributes in Unicon". doc/utr/utr28.rst. 2026.

\bibitem[12]{cite-rfc792}
J. Postel. "Internet Control Message Protocol". RFC 792. 1981.

\bibitem[13]{cite-rfc2236}
W. Fenner. "Internet Group Management Protocol, Version 2". RFC 2236. 1997.

\bibitem[14]{cite-rfc3376}
B. Cain and S. Deering and I. Kouvelas and B. Fenner and A. Thyagarajan. "Internet Group Management Protocol, Version 3". RFC 3376. 2002.

\bibitem[15]{cite-rfc1071}
R. Braden and D. Borman and C. Partridge. "Computing the Internet Checksum". RFC 1071. 1988.

\end{thebibliography}
\end{document}
